\section{Overview}

This section presents a comprehensive Failure Modes and Effects Analysis (FMEA) of the Feedback-Coupled Cryptographic Observation (RCO) protocol. The objective of this analysis is to identify, classify, and formally evaluate all potential failure modes that may arise within the system, quantify their impact on system integrity, and demonstrate that the protocol either prevents such failures or bounds their effects under well-defined assumptions.

\begin{table}[htbp]
\footnotesize
\centering
\caption{FMEA Summary of RCO Protocol Components and Failure Modes}
\label{tab:fmea-summary}

\renewcommand{\arraystretch}{1.25}

\begin{tabularx}{\textwidth}{|p{2.8cm}|X|X|X|p{2cm}|p{2cm}|p{2cm}|}
\hline
\textbf{Component} & \textbf{Failure Mode} & \textbf{Cause} & \textbf{Effect on System} & \textbf{Severity} & \textbf{Detection} & \textbf{Mitigation} \\
\hline

Telemetry Producer & Malformed or adversarial data & Compromised device or faulty input & Invalid telemetry enters pipeline & High & Schema validation, hash mismatch & Reject at ingestion stage \\

\hline

Canonicalization Engine & Non-deterministic encoding & Inconsistent serialization rules & Hash inconsistency across validators & High & Cross-validator mismatch & Enforce strict canonical schema \\

\hline

Hashing Module & Weak or broken hash function & Poor algorithm choice or implementation flaw & Collision or tampering undetected & Critical & Verification inconsistencies & Use standard cryptographic hashes (SHA-2/3) \\

\hline

Lineage Engine & Broken chain or missing links & State loss or corruption & Loss of historical integrity & Critical & Lineage verification failure & Recompute lineage, reject chain \\

\hline

State Binding Module & Incorrect state association & State desynchronization or tampering & Replay or semantic forgery & Critical & Validation failure in PoR & Enforce deterministic state hashing \\

\hline

Validator Network & Malicious validator behavior & Node compromise or collusion & Invalid data approved & Critical & Signature inconsistency & Threshold enforcement ($f < t$) \\

\hline

Aggregation Layer & Incorrect signature aggregation & Faulty logic or corrupted inputs & False acceptance or rejection & High & Verification mismatch & Deterministic aggregation + re-verification \\

\hline

Verification Engine & Incorrect validation decision & Logic error or corrupted inputs & System integrity violation & Critical & Re-validation mismatch & Strict cryptographic verification \\

\hline

Persistence Layer & Data corruption or loss & Storage failure or attack & Loss of auditability & High & Integrity checks via lineage & Append-only storage + redundancy \\

\hline

Audit \& Query Service & Incorrect query results & Incomplete lineage or corrupted data & Misleading verification output & Medium & Recompute proofs & Cross-check with stored lineage \\

\hline

\end{tabularx}
\end{table}

Unlike conventional FMEA approaches that focus on mechanical or software component failures in isolation, the RCO protocol requires a multi-layered analysis that integrates cryptographic, distributed, and feedback-coupled system dynamics. In particular, failure must be defined not only in terms of component malfunction but also in terms of violation of system invariants.

Let the system state be defined as:

\begin{equation}
\mathcal{S}_t = (D_t, S_t, L_t, \Sigma_t)
\end{equation}

A failure occurs if:

\begin{equation}
\text{Valid}(D_t) = \text{false} \quad \land \quad \text{accept}(D_t) = \text{true}
\end{equation}

or if system invariants are violated:

\begin{equation}
\exists t \text{ such that } L_t \neq \mathcal{H}(H_t \concat L_{t-1})
\end{equation}

\section{Failure Classification Model}

Failures are categorized into the following classes:

\begin{equation}
\mathcal{F} = \{\mathcal{F}_{\text{representation}}, \mathcal{F}_{\text{cryptographic}}, \mathcal{F}_{\text{distributed}}, \mathcal{F}_{\text{adversarial}}, \mathcal{F}_{\text{operational}}\}
\end{equation}

Each class corresponds to a distinct layer of the system.

\section{Failure Severity Model}

Define severity as a function:

\begin{equation}
\text{Severity}: \mathcal{F} \rightarrow \mathbb{R}^+
\end{equation}

where:

\begin{equation}
\text{Severity}(f) = \|\Delta S_t\|
\end{equation}

representing the deviation in system state caused by failure $f$.

\section{Failure Probability Model}

Let:

\begin{equation}
P(f) = \Pr[\text{failure } f \text{ occurs}]
\end{equation}

The overall system failure probability is:

\begin{equation}
P_{\text{system}} = \sum_{f \in \mathcal{F}} P(f)
\end{equation}

\section{Risk Function}

Define risk as:

\begin{equation}
R(f) = \text{Severity}(f) \cdot P(f)
\end{equation}

The protocol aims to minimize:

\begin{equation}
\sum_{f \in \mathcal{F}} R(f)
\end{equation}

\section{Failure Detection Model}

A failure is detectable if:

\begin{equation}
\exists \mathcal{D}: \mathcal{S}_t \rightarrow \{\text{true}, \text{false}\}
\end{equation}

such that:

\begin{equation}
\mathcal{D}(\mathcal{S}_t) = \text{true} \implies \text{failure detected}
\end{equation}

In RCO:

\begin{equation}
\mathcal{D} = \mathcal{I}_{\text{RCO}}
\end{equation}

\section{Failure Containment Principle}

The RCO protocol enforces:

\begin{equation}
\text{Failure} \implies \text{Reject}
\end{equation}

ensuring that failures do not propagate into accepted system state.

\section{Global Failure Invariant}

\begin{equation}
\forall t, \quad \text{accept}(D_t) \implies \text{no failure occurred at } t
\end{equation}

\section{Representation-Level Failures}

\subsection{Non-Canonical Encoding}

Failure occurs if:

\begin{equation}
D_1 \equiv D_2 \land \mathcal{S}(D_1) \neq \mathcal{S}(D_2)
\end{equation}

This leads to:

\begin{equation}
H_1 \neq H_2
\end{equation}

\subsection{Impact}

\begin{equation}
\text{Severity} = \|\Delta L_t\|
\end{equation}

\subsection{Mitigation}

RCO enforces deterministic serialization:

\begin{equation}
\forall D, \quad \mathcal{S}(D) \text{ is canonical}
\end{equation}

\section{Cryptographic Failures}

\subsection{Hash Collision}

Failure occurs if:

\begin{equation}
\exists D_1 \neq D_2 \text{ such that } \mathcal{H}(D_1) = \mathcal{H}(D_2)
\end{equation}

\subsection{Probability}

\begin{equation}
P(f) = \epsilon_{\mathcal{H}}
\end{equation}

\subsection{Impact}

\begin{equation}
\text{Severity} \approx \text{maximum}
\end{equation}

\subsection{Mitigation}

Use collision-resistant hash functions.

\subsection{Signature Forgery}

Failure occurs if:

\begin{equation}
\text{Verify}(M, \Sigma, PK) = \text{true} \text{ for forged } \Sigma
\end{equation}

\subsection{Probability}

\begin{equation}
P(f) = \epsilon_{\text{sig}}
\end{equation}

\section{Distributed System Failures}

\subsection{Validator Failure}

If validator $V_i$ fails:

\begin{equation}
\sigma_i = \emptyset
\end{equation}

System continues if:

\begin{equation}
|\mathcal{V}_{\text{active}}| \geq t
\end{equation}

\subsection{Network Delay}

Messages delayed:

\begin{equation}
T_{\text{delay}} > T_{\text{threshold}}
\end{equation}

\subsection{Impact}

\begin{equation}
\text{Liveness reduced}
\end{equation}

\section{Adversarial Failures}

\subsection{Injection Attempt}

\begin{equation}
D_t' \not\equiv D_t
\end{equation}

\subsection{Detection}

\begin{equation}
\mathcal{I}_{\text{RCO}}(D_t') = \text{reject}
\end{equation}

\section{Operational Failures}

\subsection{Storage Failure}

Loss of record:

\begin{equation}
\mathcal{R}_t \notin \mathcal{S}_{\text{log}}
\end{equation}

\subsection{Mitigation}

Replication:

\begin{equation}
\mathcal{S}_{\text{dist}} = \bigcup \mathcal{S}_i
\end{equation}

\section{Quantitative Risk Bound}

The total system risk is:

\begin{equation}
R_{\text{system}} \leq \epsilon_{\mathcal{H}} + \epsilon_{\text{sig}} + \epsilon_{\text{threshold}}
\end{equation}

\section{Conclusion}

This section establishes a formal framework for analyzing failures in the RCO protocol. By defining failure modes across multiple layers and quantifying their probabilities and impacts, we demonstrate that the protocol either prevents failures or bounds their effects under standard assumptions. This provides strong assurance of system robustness and reliability.

\section{Component-Level Failure Analysis and Propagation Modeling}

The RCO protocol is composed of multiple deterministic components, each responsible for enforcing a specific invariant. A comprehensive failure analysis requires evaluating each component independently and then analyzing how failures propagate through the system pipeline. This section presents a detailed component-level FMEA and establishes formal containment properties.

\section{Component Decomposition}

Let the system pipeline be defined as:

\begin{equation}
\mathcal{P} = \mathcal{V}_{\text{verify}} \circ \mathcal{A}_{\text{agg}} \circ \mathcal{S}_{\text{sign}} \circ \mathcal{B}_{\text{state}} \circ \mathcal{C}_{\text{lineage}} \circ \mathcal{H}_{\text{hash}} \circ \mathcal{S}_{\text{canon}}
\end{equation}

Each component introduces potential failure modes.

\section{Canonicalization Component FMEA}

\subsection{Failure Mode C1: Non-Deterministic Key Ordering}

\begin{equation}
\exists D \text{ such that } \mathcal{S}_i(D) \neq \mathcal{S}_j(D)
\end{equation}

\subsection{Effect}

\begin{equation}
H_t^{(i)} \neq H_t^{(j)} \implies L_t^{(i)} \neq L_t^{(j)}
\end{equation}

\subsection{Severity}

\begin{equation}
\text{Severity}(C1) = \|\Delta L_t\| \rightarrow \text{high}
\end{equation}

\subsection{Detection}

\begin{equation}
\mathcal{I}_{\text{RCO}} \rightarrow \text{reject}
\end{equation}

\subsection{Containment Proof}

Since signature aggregation requires identical $M_t$:

\begin{equation}
M_t^{(i)} \neq M_t^{(j)} \implies \Sigma_t \text{ invalid}
\end{equation}

Thus failure is contained.

\section{Hashing Component FMEA}

\subsection{Failure Mode H1: Hash Collision}

\begin{equation}
\mathcal{H}(B_1) = \mathcal{H}(B_2), \quad B_1 \neq B_2
\end{equation}

\subsection{Propagation}

\begin{equation}
H_t^{(1)} = H_t^{(2)} \implies L_t^{(1)} = L_t^{(2)}
\end{equation}

\subsection{Probability}

\begin{equation}
P(H1) = \epsilon_{\mathcal{H}}
\end{equation}

\subsection{Containment}

No containment at protocol level; relies on cryptographic assumption.

\section{Lineage Engine FMEA}

\subsection{Failure Mode L1: Incorrect Lineage Update}

\begin{equation}
L_t' \neq \mathcal{H}(H_t \concat L_{t-1})
\end{equation}

\subsection{Propagation}

\begin{equation}
L_t' \neq L_t \implies M_t' \neq M_t
\end{equation}

\subsection{Detection}

\begin{equation}
\text{Verify}(M_t', \Sigma_t, PK) = \text{false}
\end{equation}

\subsection{Containment Theorem}

\begin{theorem}
Any lineage inconsistency results in rejection with probability $1 - \epsilon_{\text{sig}}$.
\end{theorem}

\begin{proof}
Invalid lineage produces mismatched message $M_t'$. Without valid signatures, verification fails.
\end{proof}

\section{State Binding FMEA}

\subsection{Failure Mode S1: Incorrect State Hash}

\begin{equation}
\mathcal{H}_S(S_t') \neq \mathcal{H}_S(S_t)
\end{equation}

\subsection{Propagation}

\begin{equation}
M_t' \neq M_t
\end{equation}

\subsection{Detection}

\begin{equation}
\text{Verify}(M_t', \Sigma_t, PK) = \text{false}
\end{equation}

\section{Validator Component FMEA}

\subsection{Failure Mode V1: Malicious Validator}

\begin{equation}
\sigma_i = \text{Sign}(M_t', K_i)
\end{equation}

\subsection{Condition}

\begin{equation}
|\mathcal{C}| < t
\end{equation}

\subsection{Containment}

\begin{equation}
\text{Aggregate}(\sigma_i) \text{ invalid}
\end{equation}

\subsection{Failure Mode V2: Validator Crash}

\begin{equation}
\sigma_i = \emptyset
\end{equation}

\subsection{Impact}

\begin{equation}
\text{Liveness degraded}
\end{equation}

\section{Aggregation Layer FMEA}

\subsection{Failure Mode A1: Insufficient Signatures}

\begin{equation}
|\{\sigma_i\}| < t
\end{equation}

\subsection{Outcome}

\begin{equation}
\text{reject}
\end{equation}

\subsection{Failure Mode A2: Invalid Aggregation}

\begin{equation}
\Sigma_t \neq \text{Aggregate}(\sigma_i)
\end{equation}

\subsection{Detection}

\begin{equation}
\text{Verify}(\Sigma_t) = \text{false}
\end{equation}

\section{Verification Engine FMEA}

\subsection{Failure Mode VFY1: False Acceptance}

\begin{equation}
\text{Verify}(M_t', \Sigma_t', PK) = \text{true}
\end{equation}

\subsection{Probability}

\begin{equation}
P(VFY1) = \epsilon_{\text{sig}}
\end{equation}

\section{Failure Propagation Model}

Let failure at stage $i$ be $f_i$. The propagated failure is:

\begin{equation}
f_{i+1} = \mathcal{T}(f_i)
\end{equation}

where $\mathcal{T}$ is the transformation function.

\subsection{Propagation Bound}

\begin{equation}
\|f_{i+1}\| \leq L_i \|f_i\|
\end{equation}

\subsection{Containment Condition}

\begin{equation}
\exists j \text{ such that } \mathcal{I}_{\text{RCO}}(f_j) = \text{reject}
\end{equation}

\section{Global Containment Theorem}

\begin{theorem}
All non-cryptographic failures are detected and contained by the ingestion pipeline.
\end{theorem}

\begin{proof}
Each failure alters $M_t$. Without valid threshold signatures, verification fails. Thus, failure cannot propagate into accepted state.
\end{proof}

\section{Failure Isolation Property}

\begin{equation}
f_i \not\implies f_j \quad \text{for } i \neq j
\end{equation}

\section{Cascading Failure Prevention}

\begin{equation}
\text{Failure} \rightarrow \text{Reject} \rightarrow \text{No propagation}
\end{equation}

\section{Quantitative Propagation Bound}

\begin{equation}
\Pr[\text{propagated failure}] \leq \epsilon_{\text{sig}} + \epsilon_{\mathcal{H}}
\end{equation}

\section{Conclusion}

This section provides a detailed component-level failure analysis and demonstrates that failures are either cryptographically negligible or deterministically contained. The propagation model and containment proofs establish that the RCO protocol prevents cascading failures and maintains system integrity under a wide range of conditions.

\section{Extreme Edge Cases, Byzantine Scenarios, and Resilience Modeling}

The preceding analysis has established that individual component failures are either prevented or contained by the RCO protocol. This section extends the analysis to extreme conditions, including Byzantine adversaries, simultaneous multi-component failures, and worst-case system stress scenarios. The objective is to characterize the precise boundaries under which the system maintains its guarantees and to formally define the conditions under which degradation or collapse may occur.

\section{Extreme Edge Case Model}

We define an extreme scenario as one in which multiple failure modes occur simultaneously:

\begin{equation}
\mathcal{F}_{\text{extreme}} = \{f_1, f_2, \dots, f_k\}
\end{equation}

with $k$ maximized under system constraints.

The combined effect is:

\begin{equation}
\Delta S_t^{\text{extreme}} = \sum_{i=1}^{k} \Delta S_t^{(f_i)}
\end{equation}

The protocol must ensure:

\begin{equation}
\text{accept}(\mathcal{F}_{\text{extreme}}) = \text{false}
\end{equation}

unless all invariants hold.

\section{Byzantine Failure Model}

Let $\mathcal{C} \subseteq \mathcal{V}$ denote the set of compromised validators.

\begin{equation}
|\mathcal{C}| = f
\end{equation}

\subsection{Byzantine Capability}

The adversary may:

\begin{equation}
\mathcal{A}_{\text{Byz}} = \{\text{sign invalid messages}, \text{withhold signatures}, \text{collude}\}
\end{equation}

\subsection{Threshold Condition}

The system remains secure if:

\begin{equation}
f < t
\end{equation}

\subsection{Failure Condition}

System compromise occurs if:

\begin{equation}
f \geq t
\end{equation}

\subsection{Byzantine Collapse Theorem}

\begin{theorem}
If $f \geq t$, the adversary can produce valid signatures for arbitrary messages.
\end{theorem}

\begin{proof}
The adversary controls at least $t$ validators, enabling construction of $\Sigma_t$ for any $M_t$. Thus verification succeeds.
\end{proof}

\section{Multi-Failure Interaction Model}

Consider failures across multiple layers:

\begin{equation}
f = (f_{\text{canon}}, f_{\text{hash}}, f_{\text{lineage}}, f_{\text{sig}})
\end{equation}

\subsection{Joint Failure Probability}

\begin{equation}
P(f) = \prod_{i} P(f_i)
\end{equation}

\subsection{Upper Bound}

\begin{equation}
P(f) \leq \epsilon_{\mathcal{H}} \cdot \epsilon_{\text{sig}} \cdot \epsilon_{\text{threshold}}
\end{equation}

which is negligible.

\section{Worst-Case Scenario Analysis}

\subsection{Simultaneous Hash Collision and Signature Forgery}

\begin{equation}
\mathcal{H}(D_1) = \mathcal{H}(D_2) \land \text{Forge}(\Sigma)
\end{equation}

\subsection{Probability}

\begin{equation}
P = \epsilon_{\mathcal{H}} \cdot \epsilon_{\text{sig}}
\end{equation}

\subsection{Impact}

\begin{equation}
\text{Severity} \rightarrow \text{maximum}
\end{equation}

\section{System Collapse Conditions}

The system collapses if any of the following hold:

\begin{equation}
\text{(i)} \quad f \geq t
\end{equation}

\begin{equation}
\text{(ii)} \quad \epsilon_{\mathcal{H}} \approx 1
\end{equation}

\begin{equation}
\text{(iii)} \quad \epsilon_{\text{sig}} \approx 1
\end{equation}

\begin{equation}
\text{(iv)} \quad \mathcal{S} \text{ is non-canonical}
\end{equation}

\section{Graceful Degradation Model}

When partial failures occur:

\begin{equation}
\text{Failure} \rightarrow \text{Reject} \rightarrow \text{Retry}
\end{equation}

\subsection{Liveness Degradation}

\begin{equation}
\text{Throughput} \downarrow \quad \text{Latency} \uparrow
\end{equation}

\subsection{Safety Preservation}

\begin{equation}
\text{Invalid data never accepted}
\end{equation}

\section{Recovery Mechanisms}

\subsection{Validator Recovery}

Failed validators rejoin:

\begin{equation}
\mathcal{V}_{\text{active}} \rightarrow \mathcal{V}
\end{equation}

\subsection{State Reconstruction}

Using lineage:

\begin{equation}
L_t \rightarrow \{D_0, \dots, D_t\}
\end{equation}

\subsection{Replay Recovery}

Buffered inputs are reprocessed:

\begin{equation}
\mathcal{B} \rightarrow \mathcal{I}_{\text{RCO}}
\end{equation}

\section{Resilience Theorem}

\begin{theorem}
If $f < t$ and cryptographic assumptions hold, the system maintains safety under arbitrary failure combinations.
\end{theorem}

\begin{proof}
Failures either alter $M_t$ or prevent threshold formation. In both cases, verification fails or data is rejected.
\end{proof}

\section{Long-Term Stability}

Over time horizon $T$:

\begin{equation}
\Pr[\text{system failure over } T] \leq T \cdot (\epsilon_{\mathcal{H}} + \epsilon_{\text{sig}} + \epsilon_{\text{threshold}})
\end{equation}

\section{Adversarial Saturation Scenario}

Adversary attempts continuous attack:

\begin{equation}
\mathcal{A} = \{\mathcal{A}_1, \dots, \mathcal{A}_T\}
\end{equation}

\subsection{Bound}

\begin{equation}
\Pr[\text{success}] \leq T \cdot \epsilon
\end{equation}

\section{Global Resilience Invariant}

\begin{equation}
\forall t, \quad \text{accept}(D_t) \implies \text{system integrity preserved}
\end{equation}

\section{Safety vs Collapse Boundary}

\begin{equation}
\text{Safe Region}: f < t
\end{equation}

\begin{equation}
\text{Collapse Region}: f \geq t
\end{equation}

\section{Formal Safety Envelope}

\begin{equation}
\mathcal{E}_{\text{safe}} = \{(f, \epsilon_{\mathcal{H}}, \epsilon_{\text{sig}}) \mid f < t, \epsilon \ll 1\}
\end{equation}

\section{Conclusion}

This section establishes the behavior of the RCO protocol under extreme conditions, including Byzantine adversaries, multi-failure scenarios, and worst-case cryptographic events. The analysis demonstrates that the system maintains safety under all conditions except those that violate fundamental cryptographic assumptions or threshold constraints. Even in degraded states, the protocol preserves integrity while sacrificing liveness, ensuring that invalid telemetry never propagates into system state. This confirms the robustness of the RCO protocol as a high-assurance system suitable for safety-critical and adversarial environments.